\begin{textsample}{POS Dim 3 – rn – Score 15.00 – rn/rn\_inf\_26.txt}  \label{ex:f3_pos_059}
CAMPO DE EXPERIÊNCIAS “ \textbf{ESCUTA}, FALA, PENSAMENTO E IMAGINAÇÃO ” Direitos de aprendizagem e desenvolvimento - CONVIVER com \textbf{crianças} e \textbf{adultos} em situações comunicativas cotidianas, constituindo modos de pensar, imaginar, sentir, narrar, dialogar e conhecer. - \textbf{BRINCAR} com parlendas, trava-línguas, adivinhas, memória, rodas, \textbf{brincadeiras} cantadas, jogos e textos de imagens, escritos e outros, ampliando o repertório das manifestações culturais da tradição local e de outras culturas, enriquecendo sua linguagem oral, corporal, musical, \textbf{dramática}, escrita, dentre outras. - PARTICIPAR de rodas de conversa, de relatos de experiências, de contação e leitura de histórias e poesias, de construção de narrativas, da elaboração, descrição e representação de papéis no faz de conta, da exploração de materiais impressos e de variedades linguísticas, construindo diversas formas de organizar o pensamento. - EXPLORAR \textbf{gestos}, expressões, \textbf{sons} da língua, rimas, imagens, textos escritos, além dos sentidos das palavras, nas poesias, parlendas, canções e nos enredos de histórias, apropriando -se desses elementos para criar novas falas, enredos, histórias e escritas convencionais ou não. - EXPRESSAR sentimentos, ideias, percepções, \textbf{desejos}, necessidades, pontos de vista, informações, dúvidas e descobertas, utilizando múltiplas linguagens, considerando o que é comunicado pelos \textbf{colegas} e \textbf{adultos}. - CONHECER -SE e reconhecer suas preferências por pessoas, \textbf{brincadeiras}, lugares, histórias, autores, gêneros linguísticos, e seu interesse em produzir com a linguagem verbal. ( BRASIL, 2018 ) A comunicação é um processo que ocorre com a \textbf{criança} por meio dos diversos modos e linguagens ( movimento, olhar, postura, sorriso, choro, dentre outros ) e é participando de eventos comunicativos desde que nasce, no contato com o \textbf{adulto}, que a \textbf{criança} \textbf{escuta} e experimenta ações e reações.
A partir dos sentidos dados pelos outros, nas situações de interação, ela tenta, aos poucos, compreender o mundo que lhe é apresentado.
A linguagem verbal, portanto, associada a outras linguagens, vai representando, à medida que a \textbf{criança} cresce, a maneira encontrada para interagir com os \textbf{adultos}, com outras \textbf{crianças} e com o meio físico que lhe cerca.
Na Educação \textbf{Infantil}, essa articulação continua a ocorrer de forma contextualizada e sistematizada, através das práticas sociais que fazem parte do universo cultural vivenciado pela \textbf{criança}, sendo a aprendizagem da língua materna, \textbf{progressivamente} ampliada com a apropriação do vocabulário.
Para tanto, experiências necessitam ser planejadas com vistas à participação na cultura humana, ou seja, nas práticas sociais.
Logo, é preciso inserir a \textbf{criança} enquanto protagonista de seu processo educativo, seja verbalizando suas opiniões, ou ainda, escutando histórias, envolvendo -se em situações de diálogos, de descrições, de narrativas produzidas por elas ou com seus pares, experienciando múltiplas formas de linguagem.
Assim, a \textbf{criança} poderá experimentar seus \textbf{sons}, diferenciar modos de falar, de escrever e refletir sobre as causas de tais diferenças.
Nessa perspectiva, é papel da Educação \textbf{Infantil}, possibilitar, através da mediação do professor, experiências em que a \textbf{criança} possa fazer uso da linguagem oral e escrita, por meio das \textbf{brincadeiras} de roda, \textbf{brincadeiras} com textos, rimas, jogos cantados, formas de comunicação presentes no cotidiano, etc., a fim de construir concepções acerca da linguagem escrita, utilizando -a em situações significativas a partir de contextos reais, assim como a escrita é utilizada socialmente, sem que haja, com isso, obrigatoriedade de conclusão do processo de apropriação da língua escrita.
Sobre a aprendizagem da linguagem oral, também é relevante que as práticas planejadas considerem, nos momentos de interação, oportunidades em que as \textbf{crianças} experimentem situações de narrar, descrever, explicar, relatar, ouvir e argumentar com outros \textbf{colegas} e com os professores.
Tais práticas de leitura e de escrita, propiciadas em meio às situações \textbf{lúdicas} e ao exercício da \textbf{escuta}, contribuem para que as \textbf{crianças} exercitem sua imaginação e pensamento e com isso façam uso dessas linguagens.
É, portanto, em um ambiente linguisticamente possibilitador, que as \textbf{crianças} desenvolvem seus saberes, participando das práticas sociais significativas com essas linguagens, contextualizadas, reconhecendo seus usos e funções.
Assim, a organização do trabalho pedagógico pode se dar a partir da ludicidade, de \textbf{brincadeiras} com \textbf{sons}, terminações de palavras, rimas, cantigas ou ainda diferenciação de \textbf{sons}, modos de falar e escrever.
Além disso, é fundamental a exploração dessas linguagens usando os nomes das \textbf{crianças}, histórias que elas gostam de ouvir ou \textbf{contar}, \textbf{brincadeiras} que gostam de \textbf{brincar} ou coisas de que gostam ou se \textbf{interessam} por fazer/saber.

% matched lemmas: adulto, brincadeira, brincar, colega, contar, criança, desejo, dramático, escuta, gesto, infantil, interessar, lúdico, progressivamente, som
\end{textsample}
